% Generated draft: requires mathematical and source review.
\documentclass[11pt]{article}
\usepackage[margin=25mm]{geometry}
\usepackage{fontspec}
\setmainfont{Latin Modern Roman}
\usepackage{amsmath,amssymb,mathtools,unicode-math}
\setmathfont{Latin Modern Math}
\usepackage{enumitem,needspace,xurl,hyperref}
\setlength{\parindent}{0pt}
\setlength{\parskip}{5pt}
\setlength{\emergencystretch}{3em}
\newcommand{\N}{\mathbb{N}}
\newcommand{\Z}{\mathbb{Z}}
\newcommand{\Q}{\mathbb{Q}}
\newcommand{\R}{\mathbb{R}}
\newcommand{\C}{\mathbb{C}}
\newcommand{\F}{\mathbb{F}}
\newcommand{\A}{\mathbb{A}}
\newcommand{\PP}{\mathbb{P}}
\newcommand{\E}{\mathbb{E}}
\newcommand{\eps}{\varepsilon}
\newcommand{\dd}{\,\mathrm{d}}
\begin{document}
\section*{UnsolvedMath ID: 700006}
\textbf{Source problem code:} catalogue entry
\subsection{Definitions and mathematical statement}
A group \(G\) is right-orderable if it admits a strict total order \(<\) such that
\[
\forall g,h,x\in G,\qquad g<h\Longrightarrow gx<hx.
\]
For a specified nonzero associative ring \(R\) with identity, the group ring \(R[G]\) consists of finite formal sums \(\sum_{g\in G}r_g g\), with coefficientwise addition and multiplication determined by
\[
(rg)(sh)=(rs)(gh)\qquad(r,s\in R,\ g,h\in G).
\]
A skew field is a nonzero associative ring with identity in which every nonzero element has a multiplicative inverse. Under the unital convention, an embedding into a skew field is an injective ring homomorphism preserving the identity.

The supplied question asks whether the group ring of every right-orderable group embeds into a skew field. Once an admissible coefficient ring \(R\) is specified, its formulation is
\[
\forall G\;\Bigl(G\text{ is right-orderable}\Longrightarrow
\exists D\,\exists\iota:R[G]\hookrightarrow D\;
\bigl(D\text{ is a skew field and }\iota\text{ is a unital ring homomorphism}\bigr)\Bigr).
\]
The source does not specify the coefficient ring, its hypotheses, or its quantification. Thus this display is a formulation scheme for a specified \(R\), not a resolution of the intended coefficient domain. The unital embedding convention is also not explicit in the source.
\subsection{Short English statement}
Does the group ring of every right-orderable group admit an injective homomorphism into a skew field? The supplied statement leaves the coefficient ring and its quantification unspecified, so a fully determined mathematical statement requires clarification.
\subsection{Sources}
\begin{itemize}
\item Local UnsolvedMath registry, permanent ID \(700006\), source code \texttt{catalogue entry}, title \emph{Kourovka Notebook Problem 1.6}; supplied upstream revision \texttt{d371ee917ac9647fbb06cd6fc529deec6f3dda82}.
\item Supplied reference: \emph{The Kourovka Notebook (Unsolved Problems in Group Theory)}, original problem 1.6, first issue (1965); cited edition: 21st issue (2026), arXiv:1401.0300v48. The collection description names E. I. Khukhro and V. D. Mazurov as editors. No external source was independently consulted.
\item Attribution: the statement prefixes the question with A. I. Mal'cev; the background and proposer field name L. A. Bokut', with proposed year 1965. The supplied record does not explain the relationship between these attributions.
\item The registry labels this problem \texttt{open}. This is supplied source metadata, not an independently verified assessment of its mathematical status.
\end{itemize}
\end{document}
